\documentclass[10pt,a4paper]{article}
\usepackage[utf8]{inputenc}
\usepackage[T1]{fontenc}
\usepackage[french]{babel}
\usepackage{lmodern}
\usepackage[margin=1.8cm]{geometry}
\usepackage{booktabs,array,tabularx,longtable}
\usepackage{xcolor}
\usepackage[hidelinks]{hyperref}
\usepackage{amsmath}
\usepackage{siunitx}
\sisetup{locale=FR,detect-all}
\setlength{\parskip}{0.4em}
\setlength{\parindent}{0pt}
\newcommand{\ok}{\textcolor{green!45!black}{\textbf{OK}}}
\newcommand{\corr}{\textcolor{blue!60!black}{\textbf{Corrigé}}}
\newcommand{\ecart}{\textcolor{red!70!black}{\textbf{Écart}}}
\newcommand{\nc}{\textcolor{orange!80!black}{\textbf{Non corrigé}}}
\usepackage{xurl}
\newcommand{\vp}[1]{\path{#1}}
\newcolumntype{L}{>{\raggedright\arraybackslash}X}
\newcolumntype{P}[1]{>{\raggedright\arraybackslash}p{#1}}
\renewcommand{\arraystretch}{1.15}

\title{Validation hors calibrage : le modèle face à la littérature canine}
\date{20 novembre 2026}
\begin{document}
\maketitle

\textbf{Usage pédagogique uniquement.} Aucun résultat ne sert à décider d'un traitement réel.

\textbf{Qui a fait quoi.} Les simulations, les corrections et ce rapport ont été faits par Claude Code (Anthropic) dans une session de travail du projet. Trois sous-agents de la même session ont fait les recherches bibliographiques (hémodynamique, gaz du sang, anesthésie) via Europe PMC, PMC et des moteurs de recherche. Ce n'est \emph{pas} une relecture par des pairs ni une relecture extérieure. Chaque valeur publiée porte son niveau d'accès : \textbf{lu} (passage ou tableau du texte intégral lu), \textbf{résumé} (abstract seul), \textbf{secondaire} (chiffre lu dans un article qui cite la source).

\section*{Résumé}
Les vérifications précédentes (\vp{VERIFICATION_LITTERATURE.md}, \vp{VERIFICATION_MALADIES_MEDICAMENTS.tex}, \vp{VERIFICATION_ECG.md}, \vp{VERIFICATION_PCG.md}, calibration de l'anesthésie) comparaient surtout le modèle aux données qui avaient servi à le caler. Ici, on rejoue \emph{d'autres} protocoles publiés : hémorragie, remplissage, apnée, oxygénation sous anesthésie, réveil, hypothermie, perfusions continues, poids corporel, compensation acide-base, pneumonie, crise addisonienne, troubles du rythme.

Résultat : un défaut du moteur (convergence du cycle cardiaque), une loi physiologique fausse (compensation respiratoire) et onze écarts de calage. Tous ont été corrigés sauf quatre, qui demandent un mécanisme absent du modèle (tableau~\ref{tab:bilan}). Les corrections ont été validées par l'utilisateur le 20/11/2026, y compris la régénération du golden et des fixtures Python.

\begin{table}[h]\small
\caption{Bilan. Avant : moteur du 19/11/2026 (commit \texttt{386478e}). Après : ce rapport.}\label{tab:bilan}
\begin{tabularx}{\linewidth}{P{3.6cm}P{3.0cm}P{3.0cm}LP{1.5cm}}
\toprule
Grandeur & Avant & Après & Littérature (accès) & Statut\\
\midrule
DC en hypovolémie de 30 \% & 1,79 L/min (FC$\times$VS 1,01) & 1,53 (FC$\times$VS 1,53) & cohérence interne & \corr\\
Pente $\Delta$PaCO$_2$/$\Delta$HCO$_3^-$ & 0,41 & 0,74 & 0,74 (Madias 1984, résumé) ; 0,7 (DiBartola, lu) & \corr\\
Pneumonie (A-a 47) : PaCO$_2$ ; FC & 38,1 ; 87 & 29,9 ; 118 & 28,8 (Viitanen, lu) ; 117–130 (Goggs, Young, lus) & \corr\\
Addison, crise : HCO$_3^-$ veineux & 22,5 & 15,3 & 14,7 (12,5–19,9) (Chalifoux, lu) & \corr\\
BAV 3\textsuperscript{e} degré : FC & 57 /min & 42 & 42 $\pm$ 11 ; 43–47 (van Bavel, lu) & \corr\\
Choc sous isoflurane : décès & dès PAM $<$ 40 & après 60 min cumulées & survie à PAM 40–44 (Cardillo, lu) & \corr\\
PaCO$_2$ iso 1 MAC (cible de calage) & 44,7 après la pente & 46,2 & 46,2 (Steffey 1993) & \corr\\
PaO$_2$ sous O$_2$, anesthésié & 628–660 & 523–550 & 523–579 (Kato, Itami, Mosing, lus) & \corr\\
Sternal après iso 1 MAC, 2 h & 104 min & 16 min & 14,7–18,0 (Rastas, lu) & \corr\\
Hypothermie à 60/120/180 min & $-$1,4/$-$2,7/$-$4 °C & $-$2,0/$-$2,6/$-$2,8 & $-$2,0/$-$2,6/$-$2,9 (Redondo, résumé) & \corr\\
Épargne de MAC, dex 0,5 et 3 µg/kg/h & 50 et 73 \% & 18 et 59 \% & 18 et 59 \% (Pascoe 2006, secondaire) & \corr\\
Dex 2 µg/kg/h (hors échantillon) & 69 \% & 50 \% & 43,4 \% (Acevedo-Arcique, lu) & \corr\\
Plateau de PaCO$_2$, opioïde massif & 57 (après la pente) & 65,5 & $\le$ 67 (Bailey 1987, résumé) & \corr\\
Hémorragie, vigile : PAM à 17 / 30 mL/kg & $-22$ / $-55$ mmHg & $-10$ / $-23$ & $-15$ (Shen) / $-35 \pm 10$ (Thrasher), résumés & \corr\\
FC selon le poids (5 / 55 kg) & 124 / 68 & inchangé & écart 10,5 /min (Hezzell, résumé) & \nc\\
Réveil après iso 1,54 MAC, 2 h & extubation 44 min & 39 min & station debout 26 (Lopez, résumé) & \nc\\
FR en pneumonie & 24 & 28 & 41–61 (Young, Viitanen, lus) & \nc\\
\bottomrule
\end{tabularx}
\end{table}

\section{Méthode}
\begin{itemize}
\item \textbf{Moteur} : Rust (référence normative). Python secondaire mis à jour en miroir, qui retrouve les mêmes valeurs (golden : 1862/1862 valeurs conformes).
\item \textbf{Scénarios stationnaires} :\newline\texttt{python tests/validation\_hors\_calibrage.py -{}-cmd target/release/vetsim-cli} (poids, gradient A-a, compensation acide-base, cohérence DC = FC $\times$ VS).
\item \textbf{Scénarios dynamiques} :\newline\texttt{cargo run -{}-release -p vetsim-api -{}-example validation\_hors\_calibrage -{}- \emph{scénario}}, graine 1, chien sain de 20 kg. Scénarios : \vp{hemorragie}, \vp{bolus}, \vp{apnee}, \vp{ventilation}, \vp{pao2}, \vp{captation}, \vp{reveil}, \vp{reveil2}, \vp{temperature}, \vp{epargne}, \vp{dexcri}, \vp{dex}, \vp{propofol}, \vp{atipamezole}, \vp{convergence}. La variable \vp{VETSIM_ROOT} rejoue avec une copie modifiée des modules (balayages de calage).
\item \textbf{Halogénés} : la MAC visée est la fraction \emph{télé-expiratoire} (définition des études), tenue par un réglage du vaporisateur corrigé toutes les 15 s. Une première version réglait sur la MAC cérébrale : elle oscillait et dépassait la cible. Un décès à « 2 MAC » rapporté par cette version était un artefact de ce dépassement.
\item \textbf{Hors calibrage} : aucune des données du tableau~\ref{tab:bilan} n'avait servi à caler le modèle avant ce travail. Celles qui ont servi à recaler un paramètre sont ensuite \emph{dans} l'échantillon ; le contrôle hors échantillon restant est signalé quand il existe (dexmédétomidine : Acevedo-Arcique).
\end{itemize}

\section{Défaut du moteur : cycle cardiaque non périodique}
\vp{simulate_hemo} intégrait 11 battements depuis des volumes initiaux fixes (aorte à 100 mmHg) et mesurait le dernier. En hypovolémie, l'aorte se vidait encore : la PAM et le DC du dernier battement étaient surestimés. L'indicateur \vp{converged} (variation du VS $<$ 3 \%) existait mais n'était lu nulle part, et il répondait « vrai » à $-340$ mL malgré un biais de 4 mmHg.

\begin{center}\small
\begin{tabular}{lrrrrr}
\toprule
Déficit (mL), paramètres effectifs fixés & Battements & PAM & DC & FC$\times$VS & Dérive de PAM\\
\midrule
0 & 11 & 103,0 & 3,290 & 3,280 & $2{,}7\cdot10^{-3}$\\
0 & 80 & 102,6 & 3,275 & 3,277 & $8{,}5\cdot10^{-6}$\\
510 (30 \%) & 11 & 70,2 & 1,789 & 1,008 & $1{,}1\cdot10^{-1}$\\
510 (30 \%) & 80 & 50,2 & 1,264 & 1,265 & $3{,}2\cdot10^{-6}$\\
\bottomrule
\end{tabular}
\end{center}

\textbf{Correction} (\vp{hemo.rs}, miroir \vp{physio.py}) : 22 battements (\vp{N_BEATS}), puis doublement tant que la PAM dérive de plus de 0,5 \% entre les deux derniers (\vp{MAP_DRIFT_TOL}, au plus 176). Un premier essai gardait 11 battements avec ce seuil : Rust et Python basculaient différemment pour une dérive de 0,00500 (pneumonie 0,8), d'où 22. Coût : environ deux fois plus de calcul hémodynamique (évaluation $\approx$ 5 ms).

\textbf{Contrôle} : le script stationnaire compare DC et FC$\times$VS. Avant : $+57$ \% à $-510$ mL (échec). Après : $\le 1{,}8$ \% sur tous les cas à rythme régulier. Il ne reste que la maladie du sinus ($+3{,}9$ \%), dont les pauses rendent la moyenne FC$\times$VS non comparable.

\textbf{Conséquence sur la MMVD sévère} : au régime périodique, la POG monte (s = 0,9 : 24,1 $\to$ 26,5 mmHg ; s = 1,0 : 32,4 $\to$ 38,1) et la PaO$_2$ de la MMVD 1,0 tombe de 46 à 25 mmHg. Le calage de la MMVD avait été fait sur le moteur biaisé. Avec le seuil d'œdème déjà signalé (18 mmHg contre 20–24 dans les sources, \vp{VERIFICATION_LITTERATURE.md}), c'est un point \textbf{à arbitrer}.

\section{Physiologie stationnaire}
\subsection{Compensation respiratoire des acidoses métaboliques}
La commande ventilatoire variait en $\exp(-2{,}5(\mathrm{pH}-7{,}4))$. La pente obtenue était de 0,41 mmHg de PaCO$_2$ par mmol/L de HCO$_3^-$ (MRC et Addison). Chez le chien, la pente est de 0,74, uniforme de l'acidose sévère à l'alcalose sévère (Madias 1984, \emph{J Appl Physiol} 56:1640, résumé). La table de DiBartola donne 0,7 (lue dans Mazzaferro \& Ford 2012, tableau 1-6).

\textbf{Correction} : $\exp(0{,}05 - 6{,}0\,(\mathrm{pH}-7{,}38))$ (\vp{gas::PH_GAIN}), avec un pivot au pH du chien sain : la PaCO$_2$ de 40,2 mmHg reste inchangée. Pentes obtenues : MRC 0,73–0,75 ; Addison 0,74–0,76. La session d'anesthésie utilise la même fonction.

\textbf{Effets en chaîne, recalés} :
\begin{itemize}
\item PaCO$_2$ à 1 MAC d'halogéné : 44,7 au lieu de 46,2. Coefficients ventilatoires de l'isoflurane et du sévoflurane réajustés par \texttt{anesthesia\_calibrate -{}-fit -{}-inhalants}.
\item Plateau de PaCO$_2$ de l'opioïde massif : 57 au lieu d'environ 65 mmHg. Amplitude du terme toxique (placeholder D) portée de 0,8 à 0,9.
\end{itemize}

\subsection{Pneumonie}
Le module (D, « non calibré ») ne changeait que le gradient A-a. Sans hyperventilation ni tachycardie, la PaCO$_2$ restait à 38–40 et la FC à 87.

Valeurs publiées :
\begin{itemize}
\item Viitanen 2017 (lu, $n$ = 19), avant puis après 2 jours d'hospitalisation : PaO$_2$ 84,6 puis 69,0 ; PaCO$_2$ 30,8 puis 28,8 ; A-a 30,8 puis 46,5 ; HCO$_3^-$ 18,2–18,7 ; pH 7,41.
\item Young 2022 (lu, aspiration) : A-a 47–52 ; PaCO$_2$ 24–27 ; FC 126–130.
\item Goggs 2022 (lu) : FC 117 $\pm$ 27.
\end{itemize}

\textbf{Correction} (pneumonie 1.1.0, evidence C) : \vp{vent_drive} $\times(1+5s)$, \vp{hr} $\times(1+0{,}7s)$, \vp{cl} $+5s$. Le terme chloré représente la compensation rénale de l'alcalose respiratoire, que le moteur ne calcule pas.

\begin{center}\small
\begin{tabular}{lrrrrrr}
\toprule
& PaO$_2$ & PaCO$_2$ & A-a & HCO$_3^-$ & pH & FC\\
\midrule
Modèle $s$ = 0,53 & 78 & 31,7 & 32 & 20,2 & 7,42 & 109\\
Viitanen, $<$ 2 j & 84,6 & 30,8 & 30,8 & 18,2 & 7,41 & --\\
Modèle $s$ = 0,72 & 65 & 29,9 & 47 & 19,1 & 7,42 & 118\\
Viitanen, $>$ 2 j & 69,0 & 28,8 & 46,5 & 18,7 & 7,41 & --\\
\bottomrule
\end{tabular}
\end{center}

\textbf{Non corrigé} : FR 26–29 contre 41–61 publiées. La FR du moteur ($20\sqrt{V_A/V_{A0}}$) ne représente pas la respiration rapide et superficielle.

\subsection{Crise addisonienne}
Le HCO$_3^-$ restait à 22 mmol/L, même en crise.

Valeurs publiées, HCO$_3^-$ veineux en médiane (étendue) :
\begin{itemize}
\item Chalifoux 2023 (lu) : 14,7 (12,5–19,9) en crise ; 17,6 (13,9–21,4) hors crise ; pH veineux 7,25 et 7,30.
\item Hall 2023 (lu) : 18,4 (11,1–26,0).
\end{itemize}

\textbf{Correction} (Addison 1.1.0) : anions non mesurés \vp{xa} $+8s$ (hypoaldostéronisme, azotémie).

Résultats : en crise ($s$ = 1), HCO$_3^-$ veineux 15,3 et pH veineux 7,20 ; à $s$ = 0,5, 19,5 et 7,28. Na/K : 17–24, sous le seuil de 27 comme chez 80–85 \% des chiens (Chalifoux).

\textbf{Reste} : glycémie constante. L'hypoglycémie touche 10 à 16 \% des cas (Chalifoux ; Guzmán Ramos 2022, lu). Le module ne la tire qu'avec une probabilité de 0,5, et jamais en mode déterministe.

\subsection{Bloc AV du troisième degré}
\vp{vent_escape_rate} valait 42 /min en fréquence \emph{intrinsèque}. En boucle fermée, le baroréflexe et le gain sympathique des foyers la portaient à 57 /min. Or la fréquence publiée de 42 $\pm$ 11 /min ($n$ = 97) est une fréquence \emph{observée}, comme les 43–47 /min de van Bavel 2023 (lu, chiens conscients).

\textbf{Correction} : 27 /min intrinsèques, soit 42,2 /min observés. Le contrôle de \vp{check_rhythm.py} porte désormais sur l'évaluation complète. Il échoue avec l'ancienne valeur (57 /min, hors de 31–53).

\textbf{À noter} : DC 2,05 L/min, DFG divisé par deux et créatinine à 300 µmol/L en BAV complet. Le modèle n'a pas le remodelage chronique qui rétablit le débit jusqu'à 90 \% (Loen 2022, lu).

\subsection{Conformes, sans correction}
\begin{itemize}
\item \textbf{Fibrillation atriale} : FC 129, 160 et 204 à $s$ = 0,2, 0,5 et 0,8. Menaut 2005 (résumé) : 120 sans cardiopathie, 155 avec cardiopathie, 203 en ICC.
\item \textbf{Gradient A-a du chien sain} : 0,6 mmHg avec un QR de 0,8, contre 6,3 chez Haskins 2005 (résumé). La PaO$_2$ du modèle (98,9) colle pourtant à celle de Haskins (99,5) : l'écart vient de la convention de calcul de la PAO$_2$, pas du modèle.
\end{itemize}

\subsection{FC selon le poids (non corrigé)}
\textbf{Constat} : l'allométrie $87\,(bw/20)^{-0{,}25}$ donne 124 /min à 5 kg et 68 à 55 kg. Hezzell 2013 (résumé, 5000 chiens) ne trouve que 10,5 /min d'écart entre 5 et 55 kg ; Lamb 2010 et Ferasin 2010 (résumés) ne trouvent aucune corrélation.

\textbf{Essai retiré} : une FC indépendante du poids a été essayée. Elle donne 97,5 /min à 5 kg et 80,7 à 55 kg, mais le débit du petit chien chute : DFG 2,82 au lieu de 3,47 mL/min/kg, créatinine 93 au lieu de 76 µmol/L. C'est l'inverse de Bexfield 2008 (résumé), où le DFG/kg est plus élevé chez les petits chiens.

\textbf{Ce qu'il faudrait} : une allométrie de la taille du cœur (volume d'éjection relatif plus grand chez les petits chiens), non sourcée ici. Le modèle reste aussi sans gradient de DFG/kg selon le poids (3,47 de 3 à 60 kg).

\section{Hémodynamique dynamique}
\subsection{Hémorragie (non corrigé)}
Retraits de 5 \% du volume sanguin (85 mL/kg) toutes les 5 min, chien de 20 kg.

\begin{center}\small
\begin{tabular}{lrrrrrr}
\toprule
Perte (\% ; mL/kg) & 10 ; 8,5 & 20 ; 17 & 25 ; 21 & 30 ; 25,5 & 35 ; 30 & 40\\
\midrule
Vigile : PAM (base 102,6) & 94,3 & 80,9 & 72,9 & 64,5 & 48,2 & décès\\
Vigile : FC (base 87) & 107 & 140 & 161 & 185 & 204 & --\\
Iso 1,3 MAC : PAM (base 76,7) & 61,5 & 37,4 & 20,0 & décès & -- & --\\
\bottomrule
\end{tabular}
\end{center}

Valeurs publiées chez le chien vigile (résumés) :
\begin{itemize}
\item PAM tenue jusqu'à 18,9 $\pm$ 1,8 mL/kg (point de rupture, Thrasher 2001) ;
\item $-15 \pm 3{,}3$ mmHg à 20 mL/kg (Shen 1990) ;
\item $-23$ mmHg et FC $+89$ à 26 mL/kg (Vatner 1974) ;
\item $-35 \pm 10$ mmHg à 30 mL/kg (Thrasher 1998).
\end{itemize}

Sous isoflurane :
\begin{itemize}
\item PAM 62 $\to$ 47 ($-24$ \%) à 20 mL/kg (Seo 2025, lu) ;
\item survie jusqu'à 40 $\pm$ 12 mL/kg pour une PAM cible de 40 mmHg (Cardillo 2024, lu).
\end{itemize}

\textbf{Constat} : le modèle perd de la pression dès les premiers millilitres et chute environ deux fois trop sous isoflurane.

\textbf{Mécanismes absents} : remplissage transcapillaire, contraction splénique, vasopressine et angiotensine. Aucune valeur canine n'a été cherchée pour les caler.

\textbf{Corrigé au passage} : le décès « choc irréversible » survenait au premier pas sous PAM 40, alors que la source citée par le modèle (Morales 1999, Granot 1985) parle de 1 à 3 h vers 40 mmHg. La session tolère maintenant 60 min cumulées sous ce seuil (\vp{shock_tol_min}). Le moteur stationnaire, lui, n'a pas de notion de temps et conclut encore au décès sous 40 mmHg. Conséquences :
\begin{itemize}
\item à 2 MAC d'isoflurane en ventilation contrôlée, le chien survit (PAM 46) au lieu de mourir ;
\item sous isoflurane, l'hémorragie mortelle passe de 20 à 30 \%.
\end{itemize}

\subsection{Bolus de cristalloïde}
20 mL/kg en 15 min chez le chien vigile :
\begin{itemize}
\item FC 87 $\to$ 73 et PAM 103 $\to$ 110 ;
\item DC $+10$ \% et VS 37,7 $\to$ 49,3 mL ;
\item Ht 45 $\to$ 38,3, soit environ 74 \% du volume retenu à 15 min, 46 \% à 30 min et 17 \% à 60 min.
\end{itemize}

\textbf{Littérature} :
\begin{itemize}
\item environ 25 \% retenus à 30 min et encore $+18$ \% de volume sanguin à 4 h, après 80 mL/kg (Silverstein 2005, secondaire) ;
\item VS $+8$ à $+15$ \% après 10 puis 20 mL/kg sous isoflurane (Seo 2025, lu).
\end{itemize}

\textbf{Écart} : la constante de sortie vasculaire (\vp{crystalloid_tau_min} 30, D) donne une seule exponentielle, sans la queue lente observée. Les protocoles publiés diffèrent trop (dose, anesthésie) pour la caler : non corrigé.

\section{Anesthésie}
\subsection{Oxygénation sous O$_2$}
Le modèle n'avait pas de shunt d'atélectasie et donnait une PaO$_2$ de 628 à 660 mmHg sous FiO$_2$ $\approx$ 1.

\textbf{Correction} : un terme A-a de 110 mmHg (\vp{atelectasis_aa_mmhg}), proportionnel à la profondeur et à $(F_\mathrm{I}O_2-0{,}21)/0{,}79$, donc nul à l'air.

Résultats :
\begin{itemize}
\item isoflurane 1,1 MAC, ventilation contrôlée : 549,9 mmHg, contre 547 [495–616] (Itami 2019, lu) ;
\item sévoflurane 1,3 MAC, ventilation spontanée : 523,6 mmHg, contre 523 (513–539) (Kato 2024, lu) ;
\item Mosing 2025 (lu) : 570–579.
\end{itemize}

\subsection{Apnée}
Propofol 6 mg/kg injecté en 5 s, à l'air : SaO$_2$ artérielle $<$ 90 \% à 28 s.

McNally 2009 (résumé) mesure \emph{par oxymétrie de pouls} une SpO$_2$ $<$ 90 \% à 69,6 $\pm$ 10,6 s, mais une PaO$_2$ artérielle déjà à 62 mmHg et une SaO$_2$ à 82 \% à 30 s. Le modèle est donc cohérent avec la mesure artérielle. Il n'a pas le retard du capteur de pouls.

Après pré-oxygénation, le modèle (flow-by, F$_\mathrm{I}$O$_2$ 0,4, D) désature à 88 s. Publié : 66 $\pm$ 17 s en flow-by et 187 $\pm$ 67 s au masque (Ambros 2018, résumé) ; 298 s au masque (McNally). Le modèle n'a pas de mode « masque ».

Montée de la PaCO$_2$ en apnée sous O$_2$ : environ 4 mmHg/min. Publié : de 1,3 à 2,0 mmHg/min en oxygénation apnéique à insufflation (Yokota 2019) jusqu'à environ 8 mmHg/min, valeur que j'ai calculée sur les données de Chae 2026 (résumés). Pas de correction.

\subsection{Réveil}
La conscience valait $1-\text{indice}/\texttt{loc\_mac}$, une fonction linéaire. Le décubitus sternal demandait donc une MAC cérébrale de 0,13 et arrivait 104 min après 2 h d'isoflurane à 1 MAC.

\textbf{Correction} : $1-(\text{indice}/\texttt{loc\_mac})^3$ (\vp{awareness_hill}).

\begin{center}\small
\begin{tabular}{llrrr}
\toprule
Protocole & Source & Extubation & Tête & Sternal\\
\midrule
Iso 1,0 MAC, 2 h & modèle & 10,0 & 9,0 & 15,9\\
& Rastas 2025 (lu, $n$ = 3) & 9,3–15,0 & -- & 13,7–18,0\\
Sévo 1,45 MAC, 2 h & modèle & 8,7 & 8,2 & 11,6\\
& Lozano 2009 (résumé, MAC non précisée) & 7 (5–7) & -- & 9,5 (7,3–11,8)\\
Iso 1,54 MAC, 2 h & modèle & 39,0 & 36,5 & 47,0\\
& Lopez 2009 (résumé) & -- & -- & debout 26,3 $\pm$ 7,2\\
\bottomrule
\end{tabular}
\end{center}

\textbf{Non corrigé} : l'élimination de l'isoflurane à forte dose reste trop lente. Les coefficients de partage muscle/sang et graisse/sang sont des ordres de grandeur humains non lus (D). Défaut mineur : la tête se relève parfois avant le retour de la déglutition, car deux indices différents pilotent ces jalons.

\subsection{Hypothermie peropératoire}
Le modèle refroidissait presque linéairement vers une température d'équilibre de 28 °C : $-1{,}4$, $-2{,}7$ et environ $-4$ °C à 1, 2 et 3 h.

Publié : $-2{,}0$, $-2{,}6$ et $-2{,}9$ °C à 60, 120 et 180 min (Redondo 2012, résumé, 1525 chiens). La température chute vite, puis plafonne après 60 min (Sage 2023, résumé, 1097 chiens).

\textbf{Correction} :
\begin{itemize}
\item l'excès de pertes dû à l'anesthésie s'éteint linéairement à l'approche de 35,5 °C (\vp{vasoconstriction_c}, vasoconstriction thermorégulatrice) ;
\item \vp{loss_gain_anesth} passe de 2 à 3,5.
\end{itemize}

Résultats : $-2{,}0$, $-2{,}6$ et $-2{,}8$ °C. Avec ces pertes plus fortes, le réchauffement actif (1,5 °C/h) devenait insuffisant. Recalé à 3,5 °C/h : 38,4 °C à 3 h, contre 38,2–38,7 °C avec réchauffement après environ 2 h (Rastas 2025, lu).

\subsection{Perfusions continues et épargne de MAC}
\textbf{Dexmédétomidine} : l'épargne de MAC avait été calée en supposant les concentrations mesurées par Pascoe 2006. Or la PK du modèle donne, au même débit, des concentrations 3 à 6 fois plus hautes (1,1 contre 0,198 ng/mL ; 6,7 contre 1,9). D'où 50 et 73 \% d'épargne au lieu de 18 et 59 \%. EC50 et pente ont été recalées sur les protocoles (charge puis CRI) aux concentrations du modèle : EC50 3,0 ng/mL, pente 1,27. Résultats : 18 et 59 \% ; hors échantillon, 50 \% contre 43,4 \% (Acevedo-Arcique 2014, lu). \textbf{Reste} : l'incohérence entre la PK (Dent 2019, résumé) et les concentrations mesurées de Pascoe n'est pas résolue.

\textbf{Fentanyl} : 0,15 µg/kg/min sans charge, 24–34 \% ; Ueyama 2009 (résumé), $-35$ \% avec charge. 0,6 µg/kg/min, 52–61 \% ; Steagall 2006 (résumé, clinique), $-54$ à $-66$ \%. Simões 2016 (résumé) donne $-50$ \% à 0,15 µg/kg/min avec charge. Conforme dans l'ensemble.

\textbf{Lidocaïne} : 21, 37 et 52 \% à 50, 100 et 200 µg/kg/min. Publié : 18,7 et 43,3 \% (Valverde, données de calage) ; 27,3 \% à 100 (Acevedo-Arcique, lu) ; 29 \% à 50 (Muir 2003, résumé). Un peu haut, pas corrigé. La Cp du modèle à 200 µg/kg/min est de 6,9 mg/L, contre 4,4–4,7 chez Valverde.

\textbf{Kétamine} : 0 à 10 \% d'épargne de 10 à 100 µg/kg/min. Muir 2003 (résumé) donne $-25$ \% à 10 µg/kg/min, mais cette perfusion produit environ 0,17 µg/mL, une concentration à laquelle la relation de Solano 2006 et Pypendop 2007 (résumés, base du modèle) ne prévoit presque rien. Solano lui-même rapporte des épargnes de 11 à 39 \% à 1,07 µg/mL, que la courbe de Hill du modèle sous-estime (4 \%). Les sources ne s'accordent pas : non tranché.

\subsection{Conformes, sans correction}
\begin{itemize}
\item \textbf{Dexmédétomidine IV} : à 2 µg/kg, FC $-50$/$-40$ \% et DC $-49$/$-40$ \% à 30/60 min, PAM $+10$ \%. Moran-Muñoz 2017 (lu, sous sévoflurane, 2 µg/kg puis 2 µg/kg/h) : FC $-39$ \%, DC $-42$ \%, PAM $+10$ \% à 45 min.
\item \textbf{Propofol}, dose minimale de perte de conscience :
  \begin{itemize}
  \item sans prémédication : 5,3 mg/kg (6,16 et 8,68 jusqu'à l'intubation chez Davis 2017 et Smith 2017, résumés) ;
  \item acépromazine + méthadone : 3,6 (3,48 chez Rögels 2021 ; 4,6 $\pm$ 1,3 chez Souza 2025, lu) ;
  \item dexmédétomidine 5 µg/kg + méthadone : 2,2 (2,5 $\pm$ 0,6 chez Rastas 2025, lu ; 1,8 $\pm$ 0,6 chez Bigby 2017, résumé).
  \end{itemize}
\item \textbf{Atipamézole} : la sédation est levée en 7,8 min, mais la FC ne revient pas à 80 \% de la base en 20 min. Même constat chez Turunen 2019 (résumé) et Ko 2024 (lu).
\end{itemize}

\section{Compléments du 20/11/2026 (après la première fusion)}
\subsection{Réponse cardiaque à l'hypoxémie}
\textbf{Constat.} En abaissant la FiO$_2$, le modèle restait au repos (FC 87, PAM 103, DC 3,27) jusqu'à une PaO$_2$ de 24 mmHg (SaO$_2$ 32\,\%), puis concluait au décès (dette d'O$_2$) sans transition. Il n'avait ni chémoréflexe ni réponse sympathique à l'hypoxie. La contractilité ne baissait qu'une fois la PvO$_2$ sous 20 mmHg.

\textbf{Littérature} (chien conscient, respiration spontanée ; résumés sauf mention) :
\begin{itemize}
\item à PaO$_2$ 36 mmHg : FC $+11$/min, DC $+0{,}36$ L/min, PAM et RVS inchangées (Krebs 1999) ;
\item à PaO$_2$ 28 : FC $+50$\,\%, DC $+48$\,\%, RVS $-30$\,\%, PAM $+10$\,\%, dP/dt $+60$\,\% ; la réponse persiste après dénervation sinoaortique (Krasney \& Koehler 1977) ;
\item sous anesthésie avec ventilation isocapnique, PaO$_2$ 30 : FC $+12$ à $+28$\,\%, PAM inchangée (Hubloue 2004, lu) ;
\item hypoxie sévère : excès de lactate vers une PvO$_2$ de 20 (Cain 1965) ; Ees divisée par deux au-delà d'une SaO$_2$ de 64\,\% (Walley 1988) ;
\item asphyxie : tachycardie et hypertension, puis bradycardie, hypotension et perte du pouls à 11,4 $\pm$ 2,4 min (DeBehnke 1995).
\end{itemize}

\textbf{Correction} (moteur, \vp{engine::hypoxemic_drive}, miroir Python). Une commande $h = \mathrm{clip}\bigl((50 - \mathrm{PaO_2})/(50 - 15), 0, 1\bigr)^3$, intégrée à la boucle de décompensation, agit sur trois leviers :
\begin{itemize}
\item tonus sympathique de base \vp{s0} $+6h$ ;
\item élastance $\times(1 + 3h)$ ;
\item résistances systémiques $\times(1 - 0{,}9h)$, plancher 0,2.
\end{itemize}
Nouveaux paramètres \vp{hypox_*} de \vp{healthy_dog.json}. Le mécanisme apparaît dans \vp{outcome} sous le nom \vp{hypoxemic_drive}.

\begin{center}\small
\begin{tabular}{lrrrrrr}
\toprule
FiO$_2$ & PaO$_2$ & SaO$_2$ & FC & PAM & DC & Lactate\\
\midrule
0,21 & 99 & 96 & 87 & 103 & 3,27 & 1,2\\
0,13 & 46 & 76 & 87 & 103 & 3,28 & 1,2\\
0,115 & 36 & 64 & 95 & 104 & 3,56 & 1,3\\
0,10 & 27 & 43 & 118 & 101 & 4,38 & 1,3\\
0,09 & 20 & 27 & 172 & 74 & 5,06 & 1,9\\
0,08 & \multicolumn{6}{l}{décès (dette d'O$_2$)}\\
\bottomrule
\end{tabular}
\end{center}

\textbf{Comparaison.} Aux PaO$_2$ de l'étude :
\begin{itemize}
\item 36 mmHg : FC $+11$\,\%, DC $+9$\,\%, PAM $+2$\,\%, contre $+13$\,\%, $+11$\,\% et 0 (Krebs) ;
\item 28 mmHg : FC $+30$\,\%, DC $+29$\,\%, RVS $-22$\,\%, PAM $+1$\,\%, contre $+50$, $+48$, $-30$ et $+10$\,\% (Krasney). La réponse est trop faible, car le baroréflexe la freine.
\end{itemize}
Krebs (40 min) et Krasney (10 min) ne s'accordent pas entièrement. Une pente plus raide pour atteindre Krasney faisait exploser la FC sous 22 mmHg.

\textbf{Asphyxie} (session, isoflurane 1,3 MAC, apnée sous O$_2$, scénario \vp{asphyxie}) :
\begin{itemize}
\item la PAM monte de 76 à 92 mmHg pendant la phase hypercapnique ;
\item après la désaturation (5,5 min), FC et PAM baissent progressivement (FC 105 $\to$ 73, PAM 92 $\to$ 61) ;
\item arrêt hypoxique à 10,5 min.
\end{itemize}
Il a fallu une seconde correction. La dette d'O$_2$ abaissait la FC intrinsèque, mais le sympathique compensait et la FC remontait de 86 à 132/min avant l'arrêt. Un myocarde hypoxique ne répond plus aux catécholamines : la dette bloque désormais la réponse sympathique (\vp{beta_block} $+0{,}9\cdot$DEBT, \vp{hypoxia_effects}).

\textbf{Reste} :
\begin{itemize}
\item pas de tachycardie franche au moment de la désaturation, car la dette d'O$_2$ s'accumule aussitôt ;
\item léger rebond de FC (73 $\to$ 78) dans la dernière minute ;
\item la bradycardie chémoréflexe primaire décrite à ventilation fixée (Kontos 1965) n'est pas modélisée.
\end{itemize}

\subsection{MMVD traitée (furosémide, pimobendane)}
Régimes oraux de 7 à 84 jours (simulateur de traitement du Python secondaire ; le Rust concorde sur les cas du golden).

\textbf{Potassium.} Il baissait deux fois trop :
\begin{itemize}
\item modèle avant : $-9$\,\% à J84 et $-18$\,\% à J1 à 4 mg/kg/j ;
\item publié : 4,50 en stade C (3,5 mg/kg/j) contre 4,70 en B2, soit $-4$\,\% (Adin 2020, lu) ; 4,60 $\to$ 4,50 à J84 sous 2,78 mg/kg/j (TEST, lu) ; 4,4 $\to$ 4,0 au pic puis 4,4 vers J12 (Giorgi 2022, lu).
\end{itemize}
Correction : coefficient de perte de 0,00585 à 0,0030. Résultats : $-4{,}6$\,\% à J84 et $-9{,}5$\,\% à J1, contre $-4$ (Adin) et $-9$\,\% (Giorgi).

\textbf{Créatinine.} Elle restait à 75 µmol/L sous furosémide : le DFG ne dépendait que de la PAM et du débit.
\begin{itemize}
\item publié : $+22$\,\% (0,9 $\to$ 1,1 mg/dL) après environ 12 jours à 4,1 mg/kg/j (Giorgi 2022, lu) ; BUN 20 $\to$ 25 ;
\item correction : la contraction volémique abaisse le DFG, \vp{nephron_mass} $\times\max(0{,}5,\ 1 - 3{,}4\,D/V_b)$ (furosémide 1.4.0, miroir Python) ;
\item résultats : $+16$, $+21$ et $+27$\,\% à 2,8, 4 et 8 mg/kg/j.
\end{itemize}

\textbf{Reste.} En stade C, Adin trouve une créatinine $+47$\,\% plus élevée qu'en B2, et Nicolle 2007 (résumé) un DFG de 1,7 contre 3,1 mL/min/kg en classe NYHA III–IV. Cette part d'origine cardiaque (syndrome cardiorénal) n'est pas modélisée.

\textbf{Pimobendane} : conforme sans correction. EPIC (lu) : FC inchangée à J35, fraction de raccourcissement $+2$ points. Le modèle donne DC $+0{,}4$ L/min sans hausse de FC.

\textbf{Gaz du sang et FC sous traitement.} Avant/après traitement de l'œdème cardiogénique, aucune PaO$_2$ ou SpO$_2$ chiffrée n'a été trouvée. Le modèle à 20 kg en MMVD 1,0 passe de 25 à 71 mmHg sous furosémide 4 mg/kg q8h + pimobendane. La FC publiée en ICC traitée (132–144/min, Giorgi, TEST, Yin 2026, lus) est comparable à celle du modèle (133–160).

\subsection{Axe cardio-rénal}
\textbf{Constat.} Les deux organes ne communiquaient presque pas.
\begin{itemize}
\item Le DFG ne dépendait que de la PAM (en dessous de 75 mmHg) et du débit. En MMVD terminale, malgré une PVC de 10,6 mmHg et un tonus sympathique triplé, le DFG restait à 3,50 mL/min/kg.
\item En MRC, le cœur ignorait l'anémie : à Ht 29\,\%, FC, PAM et débit étaient ceux du chien sain.
\item En mode déterministe, la PAM de la MRC baissait au lieu de monter.
\end{itemize}

\textbf{Littérature} (lue sauf mention) :
\begin{itemize}
\item DFG en classe NYHA III–IV : 1,7 contre 3,1 mL/min/kg en classe I–II (Nicolle 2007, résumé).
\item Chiens B2 non traités : DFG 2,8 mL/min/kg (Kaplan 2022).
\item Créatinine 1,5 chez les chiens traités contre 1,0 mg/dL chez les non-traités ; 32\,\% contre 1\,\% de MRC (Martinelli 2016).
\item Créatinine en stade C : $+47$\,\% par rapport au stade B2 (Adin 2020).
\item Pression veineuse rénale portée à 15 puis 30 mmHg : DFG et débit rénal inchangés (Fiksen-Olsen 1991–1992, résumés).
\item Hémodilution : débit cardiaque $+63$ à $+65$\,\% à Ht 20\,\% (Rooney 1993, résumé) ; à Hb 5 g/dL, IC $+53$\,\%, RVS $-53$\,\%, PAM $-26$\,\% (Champion 2011).
\item Pression systolique en MRC : 157 et 154 mmHg en IRIS II et III, contre 129 chez le sain (Pelander 2017) ; 156 contre 138 (Sharkey 2009, résumé).
\end{itemize}

\textbf{Mécanismes ajoutés} (Rust, miroir Python ; liens dans le schéma de l'application web) :
\begin{itemize}
\item \textbf{Cœur $\to$ rein, voie neurohormonale.} DFG $\times\max(0{,}3;\ 1 - 1{,}4\,(s - 0{,}2))$, avec $s$ le tonus sympathique (substitut du SNS et du RAAS). Résultats sans diurétique : MMVD 0,5 : 3,23 mL/min/kg ; MMVD 0,8 : créatinine $+23$\,\% ; MMVD 1,0 : DFG 1,43. Avec le furosémide ($\approx +21$\,\%), on retrouve le $+47$\,\% d'Adin en stade C.
\item \textbf{Cœur $\to$ rein, congestion veineuse.} Le terme existe (\vp{renal_k_cong}) mais reste à 0 : la littérature canine ne montre aucun effet jusqu'à 30 mmHg.
\item \textbf{Rein $\to$ cœur, anémie.} RVS $\times(\mathrm{Hb}/15)^{1{,}0}$, car la viscosité baisse. Résultats : Ht 20\,\% $\to$ DC $+56$\,\% ; Ht 15\,\% $\to$ DC $+77$\,\%, RVS $-59$\,\%, PAM $-24$\,\%. Écart : la hausse du débit passe surtout par la FC ($+69$\,\%, contre $+24$\,\% chez Champion).
\item \textbf{Rein $\to$ cœur, hypertension rénale.} Consigne du baroréflexe $+30\cdot s$ mmHg (MRC 1.2.0). Une simple hausse des RVS était annulée par le baroréflexe. Résultats : PAM 111 et 114 mmHg aux sévérités 0,5 et 0,8 ; PAS $+9$ à $+13$ mmHg. Encore sous Pelander ($+25$ mmHg). L'hypertension sévère aléatoire du module est conservée.
\end{itemize}

\textbf{Recalage de la MRC.} Le tonus sympathique monte désormais en MRC (anémie, hypertension), ce qui abaisse le DFG. Le coefficient néphronique passe de 2,4 à 2,0 pour garder les créatininémies par stade : 118, 243, 539 et 991 µmol/L aux sévérités 0,2, 0,5, 0,8 et 1,0, contre 121, 249, 512 et 828 avant. Seul le stade terminal monte.

\textbf{Effet de bord.} La crise addisonienne, avec hypovolémie et sympathique élevé, voit sa créatinine passer de 116 à 265 µmol/L (sévérité 0,8). Les publications donnent des médianes de 1,4 à 1,7 mg/dL, soit 124 à 150 µmol/L, avec 75\,\% d'azotémie (Hall 2023, Chalifoux 2023). Le modèle est donc plutôt haut : \textbf{à surveiller}.

\textbf{Non modélisé.}
\begin{itemize}
\item Expansion du volume plasmatique en IRIS III (Pelander : facteur 63,7 $\to$ 139,9 mL/kg, mesure à confirmer).
\item Hypertrophie ventriculaire gauche.
\item Troponine et NT-proBNP.
\end{itemize}

\section{Compléments du 21/11/2026}
\subsection{Seuil d'œdème pulmonaire}
Le gradient A-a s'ouvrait au-delà d'une POG de 18 mmHg, valeur écrite en dur dans \vp{gas.rs}. Les sources situent l'œdème au-delà d'environ 24 mmHg chez le chien (Guyton \& Lindsey 1959, cité par Yao \& Aoki 2025, citation secondaire lue ; écart déjà signalé dans \vp{VERIFICATION_LITTERATURE.md}).

\textbf{Correction} : seuil à 24 mmHg, constante unique \vp{engine::LA_EDEMA} (miroir Python).

\textbf{Résultats en MMVD} :
\begin{itemize}
\item sévérité 0,9 : POG 26,5 mmHg, PaO$_2$ 89 mmHg ;
\item sévérité 1,0 : POG 38,1 mmHg, PaO$_2$ 46 mmHg, contre 25 avant.
\end{itemize}

\subsection{Fréquence respiratoire en pneumonie}
\textbf{Constat} : la FR du moteur, $20\sqrt{V_A/V_{A0}}$, ne montait qu'à 28/min.

\textbf{Correction} :
\begin{itemize}
\item multiplicateur \vp{rr_mult}, qui représente la respiration rapide et superficielle ;
\item dans la pneumonie 1.2.0, $\times(1 + 1{,}1\,s)$.
\end{itemize}

\textbf{Résultats} : 42/min à $s$ = 0,53 (Viitanen : 43) et 49/min à $s$ = 0,72 (Young : 41–50). La PaCO$_2$ ne change pas.

\subsection{Créatinine de la crise addisonienne}
La hausse de créatinine apportée par l'axe cardio-rénal (section précédente) ne demande pas de correction.
\begin{itemize}
\item À sévérité 0,5 (forme modérée), la créatinine vaut 132 µmol/L, dans les médianes publiées : 1,4 à 1,7 mg/dL, soit 124 à 150 µmol/L (Hall 2023, lu).
\item Seules les crises extrêmes montent plus haut : 265 µmol/L à 0,8 et 495 à 1,0.
\end{itemize}

\subsection{Hémorragie : compensations ajoutées}
La session d'anesthésie ne compensait une hémorragie que par le baroréflexe : la PAM chutait dès les premiers mL/kg (section 4). Trois mécanismes sont ajoutés (\vp{anesthesia_model.json} 1.8.0).
\begin{itemize}
\item \textbf{Remplissage transcapillaire} vers 60\,\% du volume perdu, avec $\tau$ = 2,9 h. Ces deux valeurs reproduisent 30\,\% à 2 h (Joles 1989) et 56\,\% à 8 h (Miller \& Dale 1978).
\item \textbf{Contraction splénique} : jusqu'à 6,8 mL/kg (8\,\% du volume sanguin, Guntheroth \& Mullins 1963), commandée par le tonus sympathique. Le sang splénique est concentré, ce qui fait monter l'hématocrite.
\item \textbf{Vasopressine et angiotensine II} : RVS $\times(1 + 0{,}6\,x^4/(x^4 + 0{,}25^4))$, avec $x$ la fraction du volume perdue. La mi-hauteur suit la flambée de vasopressine entre 20 et 30 mL/kg (Wang 1988).
\end{itemize}

Résultats chez le chien conscient (retraits de 5\,\% toutes les 5 min) :
\begin{itemize}
\item 17 mL/kg : PAM $-10$ mmHg, contre $-22$ avant (Shen 1990 : $-15 \pm 3$ à 20 mL/kg) ;
\item 25\,\% : DC $-41$\,\% (Francis 2004 : $-40$\,\%) ;
\item 30 mL/kg : PAM $-23$ mmHg (Thrasher 1998 : $-35 \pm 10$).
\end{itemize}

Sous isoflurane 1,3 MAC, à 20\,\% : $-28$\,\% (Seo 2025 : $-24$\,\%).

Sources : résumés, sauf Seo (lu).

\textbf{Reste} :
\begin{itemize}
\item sous isoflurane, le décès survient vers 30 mL/kg, alors que les chiens de Cardillo 2024 tolèrent 40 $\pm$ 12 mL/kg ;
\item il manque la phase sympatho-inhibitrice (bradycardie vers 30\,\%).
\end{itemize}
Test : \vp{hemorrhage_matches_literature}, avec un contrôle d'échec sans les trois compensations.

\subsection{Banc de validation dans la CI}
\vp{crates/vetsim-core/tests/litterature.rs} reprend les comparaisons de ce rapport sous forme de tests (\vp{cargo test}). Une régression du modèle fait donc échouer la CI.
\begin{itemize}
\item \textbf{21 vérifications stationnaires} et un test d'hémorragie, chacune avec sa fourchette publiée et sa source : chien sain, pente de compensation, MRC, pneumonie, Addison, BAV, FA, axe cardio-rénal, œdème, hypoxémie, anémie, cohérence du cycle.
\item \textbf{Un scénario d'anesthésie} : PaO$_2$ sous O$_2$ et délai de l'arrêt en asphyxie.
\item \textbf{Trois contrôles d'échec}. Sans l'axe cardio-rénal, sans la viscosité de l'anémie ou sans la réponse à l'hypoxémie, la vérification correspondante doit échouer, et elle échoue.
\end{itemize}

\section{Rein dynamique (21/11/2026)}
Le rein était stationnaire : DFG en produit de facteurs, créatinine et urée à l'équilibre, sodium, eau et potassium sans bilan. Le module \vp{crates/vetsim-core/src/renal.rs}, avec ses constantes dans \vp{modules/renal/renal_model.json} (1.3.3), ajoute :
\begin{itemize}
\item le bilan du sodium, donc la variation du volume extracellulaire ($V$), iso-osmotique ; le plasma en reçoit la part $V_p/\mathrm{VEC}$ (hémo- et protéinoconcentration) ;
\item le bilan du potassium en deux compartiments, extracellulaire et intracellulaire ;
\item l'aldostérone en deux composantes, portée par la rénine ($A$) et par la kaliémie ($A_K$) ;
\item la cinétique de la créatinine et de l'urée ;
\item (1.3.0) le freinage tubulaire, la polyurie chronique d'eau libre et le rétrocontrôle tubulo-glomérulaire du DFG ;
\item (1.3.1) la libération de rénine par la macula densa et la latence de l'action natriférique de l'aldostérone.
\end{itemize}
Il ne sert que dans les évolutions temporelles, via le message \vp{renal_course} de l'API. L'état stationnaire (\vp{eval}, golden) est son point fixe : les apports de Na et de K (ceux d'un chien sain de même poids) valent l'excrétion du patient de départ, dont les FE de base s'adaptent. Le test \vp{renal_fixed_point_is_the_steady_state} le vérifie.

Équations ($b$ : état de départ, $E$ : effet du diurétique de l'anse issu de sa pharmacodynamie) :
\[
\mathrm{FE_{Na}} = \bigl(\mathrm{FE_{Na,b}} + f_\mathrm{anse}E\,(1-\varphi B)\bigr)\,A_n^{-a_\mathrm{Na}}\Bigl(\tfrac{\mathrm{PAM}}{\mathrm{PAM}_b}\Bigr)^{n_p},
\qquad
\mathrm{DFG} = \mathrm{DFG_{ss}}\Bigl(\tfrac{\mathrm{DFG}_b}{\mathrm{DFG_{ss}}}\Bigr)^{E}\,e^{-g_\mathrm{DFG}\max(0,-V)(1-E)/\mathrm{VEC}},
\]
\[
\text{diurèse} = u_0 + 1000\,\frac{U_\mathrm{Na} - \mathrm{Na_{in}}}{\mathrm{Na}_u} + w_\mathrm{max}\,W\,\mathrm{bw},\qquad
\frac{dB}{dt} = \frac{E - B}{\tau_\mathrm{on}\ \text{ou}\ \tau_\mathrm{off}},\qquad
\frac{dW}{dt} = \frac{E - W}{\tau_w},
\]
\[
U_K = \mathrm{DFG}\,\bigl(K_b\,\mathrm{FE_{K,b}} + \sigma\,(K - K_b)\bigr)\Bigl(\tfrac{\mathrm{FE_{Na}}}{\mathrm{FE_{Na,b}}}\Bigr)^{k_\mathrm{flow}} (A\,A_K)^{a_K},
\]
\[
\frac{dV}{dt} = 1000\,\frac{\mathrm{Na_{in}} - \mathrm{DFG}\cdot\mathrm{Na}\cdot\mathrm{FE_{Na}}}{\mathrm{Na}},\qquad
\frac{dA}{dt} = \frac{e^{g\,\max(0,-V)/\mathrm{VEC}}(1+g_\mathrm{md}E)(1-\mathrm{IEC}) - A}{\tau_A},\qquad \frac{dA_n}{dt} = \frac{A - A_n}{\tau_n},\qquad
\frac{dA_K}{dt} = \frac{\frac{K - k_0}{K_b - k_0} - A_K}{\tau_A},
\]
\[
\frac{dK_e}{dt} = K_\mathrm{in} + K_\mathrm{abs} - U_K - J,\qquad
\frac{dK_i}{dt} = J = \frac{B\cdot \mathrm{bw}\cdot\Delta K - K_i}{\tau_{ic}},\qquad \Delta K = K_e/\mathrm{VEC}.
\]
$\mathrm{FE_{Na,b}}$ et $\mathrm{FE_{K,b}}$ : FE de base du patient (apport $/$ charge filtrée). $\sigma$ : réserve d'excrétion au-dessus ($\sigma_\uparrow$) ou au-dessous ($\sigma_\downarrow$ = 0,13, Young 1982) de la kaliémie de base, proportionnelle au DFG (Bourgoignie 1981). $A_K$ : composante kaliémique de l'aldostérone, nulle à $k_0$ = 3,29 mmol/L (Young 1984) ; elle ne sécrète que le K (« paradoxe de l'aldostérone »). $a_K$ = 0,7 (Young 1983 ; recalé à 0,49 en 1.3.3, voir plus bas). $K_\mathrm{abs}$ : absorption d'une charge orale de K. Tous ces chiffres canins sont tirés de résumés.

La physiologie (\vp{evaluate}) est recalculée à chaque pas (0,25 h), mais seulement si un paramètre a bougé de plus de 0,2\,\% (K : 0,03 mmol/L). Une évaluation de MMVD coûte environ 200 ms.

\subsection{Étape 1 : autorégulation de la filtration}
Le premier calage a montré un défaut du rein \emph{stationnaire}. Une perte de 165 mL chez un chien de 23 kg (8\,\% du volume sanguin) abaissait le DFG de 35\,\%. Le facteur de débit (rampe de 0,5 à 1,0 fois le débit de référence) en portait la moitié, sans aucune autorégulation. Après une perfusion de furosémide, la créatinine montait de 60\,\% à 24 h, alors qu'Adin 2021 mesure $+23{,}9 \pm 12{,}4$\,\% (lu). Chez le chien en constriction cave, le DFG est conservé malgré la baisse du débit, avec une fraction de filtration accrue (résumés de recherche, plusieurs études des années 1960 à 1980).

Correction : nouveau paramètre \vp{renal_co_auto} = 0,8. Le débit ne limite la filtration qu'au-dessous de 0,8 fois le débit de référence ; à 1,0, on retrouve la rampe historique. La valeur a été choisie par balayage (1,0 / 0,8 / 0,7) sur la créatinine d'Adin à 24 h (evidence C). La MMVD n'est pas touchée : son DFG ne dépend que du facteur sympathique, et son débit reste au-dessus de la référence.

\subsection{Calage et validation}
Sources de base, accès lu : FE$_\mathrm{Na}$ 0,20\,\% et FE$_\mathrm{K}$ 9,47\,\% chez le chien sain (Troìa 2018).

\begin{table}[h]\small
\caption{Rein dynamique 1.3.3 : calage conjoint (18 constantes, 31 cibles pondérées) et validations (modèle à \texttt{dt} 0,25 h ; chronique à 1 h, Lantis à 0,5 h).}
\begin{tabularx}{\linewidth}{P{4.6cm}P{2.4cm}P{3.3cm}LP{1.3cm}}
\toprule
Protocole & Modèle & Publié & Rôle & Statut\\
\midrule
Koh 2021 IV, 2 mg/kg, $\approx$ 11,5 kg : diurèse 1 h ; 8 h & 13,9 ; 24,4 mL/kg & 12,8 ; 24,65 (lu) & calage & \ok\\
Koh 2021 IV : natriurèse 1 h ; 8 h & 1,98 ; 2,40 mmol/kg & 1,54 ; 2,62 (lu, $\Sigma$ débit $\times$ durée) & calage & \ecart\ ($+29$\,\%) ; \ok\\
Koh 2021 IV : kaliurèse 8 h rapportée à la base & $\times 1{,}80$ & $\times 2{,}27$ (lu ; chiens à jeun) & calage & \ecart\ ($-21$\,\%)\\
Koh 2021 IV : $\Delta$K à 2 h / 8 h & $-0{,}47$ / $-0{,}32$ & $-0{,}51$ / $-0{,}33$ & calage & \ok\\
Koh 2021 per os : diurèse 2 h ; 8 h & 18,5 ; 31,0 mL/kg & 18,9 ; 30,9 (lu) & calage & \ok\\
Koh 2021 per os : natriurèse 2 h ; 8 h & 2,48 ; 2,86 mmol/kg & 2,61 ; 3,55 (lu) & calage & \ok\ ; \ecart\ ($-19$\,\%)\\
Koh 2021 per os : $\Delta$K à 2 h / 8 h & $-0{,}65$ / $-0{,}48$ & $-0{,}54$ / $-0{,}46$ (lu) & calage & \ok\ (1 ET)\\
Bourgoignie 1981, chien normal, 50 mEq de KCl : pic de K ; part excrétée en 5 h & $+1{,}03$ ; 65\,\% & $+0{,}9$ à $1{,}2$ ; 61–67\,\% (résumé) & calage & \ok\\
Adin 2021, perfusion 5 h, 23,4 kg : créatinine à 24 h & $+27$\,\% & $+23{,}9 \pm 12{,}4$\,\% (lu) & calage & \ok\\
Lantis 2011, 2 mg/kg PO q12h, 30 kg : aldostérone urinaire J5 / J10 & $\times 2{,}49$ / $\times 2{,}38$ & $\times 2{,}41$ / $\times 2{,}57$ (lu) & calage & \ok\\
Lantis 2011 : $\Delta$K J5 / J10 ; poids (VEC) J5 / J10 & $-0{,}43$ / $-0{,}41$ ; $-3{,}0$ / $-2{,}7$\,\% & $-0{,}5$ / $-0{,}3$ ; $-4{,}0$ / $-1{,}7$\,\% (lu) & calage & \ok\\
Pelligand 2020, 1 à 8 mg/kg/j per os, 10 j : diurèse J9 & 25 / 40 / 75 / 95 mL/kg/j & 20 / 34 / 82 / 100 (lu) & calage & \ok\\
Pelligand 2020 : natriurèse J1 à 8 mg/kg/j ; J3 à 2 et 8 mg/kg/j & 5,05 ; 1,80 et 1,54 & 5,30 ; 1,51, base (lu) & calage & \ok\\
\midrule
Adin 2021 : créatinine pendant la perfusion (5 h) & $+0{,}3$\,\% & inchangée (lu) & validation & \ok\\
Bourgoignie 1981, rein restant (DFG $\times 0{,}27$) : part excrétée en 5 h & 33\,\% & 30–37\,\% & validation & \ok\\
Bourgoignie 1981, rein restant : pic de K & $+1{,}44$ & $+2{,}2$ à $2{,}5$ & validation & \ecart\\
Young 1976 : la régulation du K s'appuie sur l'aldostérone & charge $\times 1{,}5$ sur 7 j : $+0{,}09$ & $+8$\,\% contre $+39$\,\% sans aldostérone (résumé) & validation qualitative & \ok\\
Giorgi 2022, $\approx$ 4,1 mg/kg/j, MMVD 0,8 : créatinine J12 ; K & $+29$\,\% ; $-0{,}31$ & $+22$\,\% ; retour à la base (lu) & validation & \ok\ ; \ecart\\
Lantis 2011 : créatinine J10 & $+37$\,\% & non mesurée ; phosphorémie revenue à la base à J10 (lu) & validation qualitative & \ecart\ probable\\
Wong 1977, 1 mg/kg/j IV, 7 j : $\Delta$K & $+0{,}07$ & $-0{,}39$ (accès moteur) & validation & \ecart\\
\bottomrule
\end{tabularx}
\end{table}

\textbf{Effet de \vp{renal_co_auto} sur l'état stationnaire.} Le golden a été régénéré depuis le Python, qui reçoit le même paramètre. 41 valeurs sur 1862 sortent de la tolérance de l'ancien golden, toutes des créatinines, urées et DFG d'états à bas débit (Addison, FA, BAV, maladie du sinus, bêtabloquant, opioïde, furosémide $D$/$U$). Exemples :
\begin{itemize}
\item BAV 0,8 (débit à 62\,\% de la référence) : créatinine 668 $\to$ 258 µmol/L ;
\item Addison 0,8 : 265 $\to$ 171 µmol/L, soit 1,9 mg/dL (Hall 2023 : médiane 1,7 mg/dL, lu) ;
\item furosémide $D$ = 150 mL : 190 $\to$ 125 µmol/L.
\end{itemize}
La MMVD est inchangée, donc les calages stationnaires de Giorgi et d'Adin 2020 aussi. Rust contre le nouveau golden : 1862/1862 valeurs conformes. Les fixtures Python M1 à M4 et \vp{calibration_run_python.json} ont été régénérées.

\subsection{Chronique (1.3.0) : freinage, eau libre et rétrocontrôle du DFG}
La 1.2.0 ne valait que pour l'aigu. En prises répétées, la déplétion sodée était surestimée (VEC $-76$ mL/kg à J5 chez le chien sain, $-114$ mL/kg à J12 dans la MMVD), et la créatinine de Lantis montait de $+132$\,\%.

\textbf{Données de calage} (Pelligand 2020, lu) : beagles sains de 10,6 kg, furosémide per os en deux prises par jour, 1, 2, 5 ou 8 mg/kg/j pendant 10 jours. Natriurèse de base 1,51 mmol/kg/j et diurèse de base 0,83 mL/kg/h, qui remplacent les placeholders précédents (1 mL/kg/h) ; 5,30 mmol/kg le premier jour à 8 mg/kg/j ; retour près de la base dès J2–J3. La diurèse reste haute à J9 (20, 34, 82, 100 mL/kg/j) : la polyurie chronique est de l'eau libre, pas du sodium.

\textbf{Trois mécanismes ajoutés.}
\begin{itemize}
\item \textbf{Freinage tubulaire} $B$ : il s'installe pendant l'exposition ($\tau_\mathrm{on}$ 13,8 h en 1.3.1) et s'efface en quelques jours ($\tau_\mathrm{off}$ 202 h). Il retire jusqu'à la part $\varphi$ = 0,92 de l'effet de l'anse (hypertrophie distale, rétention post-diurétique). Evidence C.
\item \textbf{Perte du pouvoir de concentration} $W$ : elle suit l'exposition moyenne ($\tau_w$ 42,6 h en 1.3.1) et ajoute une diurèse d'eau libre ($w_\mathrm{max}$ 10,5 mL/kg/h), compensée par la boisson. Un premier essai, qui couplait l'eau libre à l'effet instantané $E$, gonflait la diurèse aiguë de Koh. Le calage avait alors divergé vers une natriurie de 600 mmol/L. Evidence C.
\item \textbf{Rétrocontrôle tubulo-glomérulaire} : pendant l'exposition, le blocage du NKCC2 de la macula densa préserve le DFG ; ensuite, le DFG baisse si le VEC reste déplété ($g_\mathrm{DFG}$ = 1,1 en 1.3.0, 1,96 en 1.3.1). Ce mécanisme répond à un fait d'Adin 2021 (résumé) : la créatinine est inchangée pendant la perfusion et ne monte qu'à 24 h, alors que le poids n'est pas récupéré. Sans lui, le modèle faisait l'inverse ($+14$\,\% à 5 h, $+8$\,\% à 24 h). Evidence C, calé sur Adin.
\end{itemize}

\textbf{Pourquoi ce troisième mécanisme.} Sans lui, Pelligand et Adin s'opposaient. Caler la natriurèse de J1 et le retour à l'équilibre demande un freinage et une rétention sodée forts ($a_\mathrm{Na}$ monte à 1,95). Ce freinage limite la déplétion après la perfusion d'Adin, et la créatinine à 24 h tombait à $+8$\,\%. Deux autres pistes ont été essayées, sans résultat :
\begin{itemize}
\item aldostérone limitée au néphron distal : la rétention entre les prises ne compense plus la fuite, et la déplétion chronique double ;
\item libération directe de rénine par la macula densa : elle aggrave la rétention sans corriger Adin.
\end{itemize}
Une fois le DFG relié à la déplétion après l'exposition, une déplétion modérée suffit pour la créatinine d'Adin.

\textbf{Résultats (1.3.0).} Lantis (10 j) : VEC $-28$ mL/kg, créatinine $+23$\,\% (contre $-76$ mL/kg et $+132$\,\% en 1.2.0). Giorgi (MMVD 0,8, J12) : créatinine $+17$\,\% contre $+22$\,\% publiés, K revenu à la base.

\textbf{Écarts restants.}
\begin{itemize}
\item \textbf{Aldostérone de Lantis} et \textbf{natriurèse aiguë de Koh} : écarts de la 1.3.0, réduits en 1.3.1 (voir ci-dessous).
\item \textbf{Hétérogénéité des sources} : chez Pelligand, deux prises de 4 mg/kg donnent 5,3 mmol/kg en 24 h ; chez Koh, une prise de 2 mg/kg en donne 4,56 en 8 h. L'apport sodé diffère : environ 2,7 mmol/kg/j dans le régime d'Adin (112 mg/100 kcal, lu), contre 1,51 mmol/kg/j mesurés chez Pelligand. Chez l'homme, la restriction sodée réduit la réponse au furosémide et accroît le freinage (résumé de recherche). Ce mécanisme n'est pas représenté, faute de données canines.
\item la natrémie basse d'Adin à 24 h n'est pas représentée (eau libre supposée exactement compensée).
\end{itemize}

\textbf{Domaine validé} : furosémide, une prise ou des prises répétées, jusqu'à 12 jours, sans levier alimentaire, charge de K ni IEC. Hors de ce domaine, l'API renvoie \vp{validated} = false et un avertissement. La voie stationnaire du diurétique (\vp{pkpd::timecourse}, $D$/$U$) reste disponible pour le golden et les calages stationnaires.

\subsection{Rénine de la macula densa et latence de l'aldostérone (1.3.1)}
Dans la 1.3.0, l'aldostérone portée par la seule déplétion montait en quelques heures et agissait aussitôt sur le Na. Elle freinait donc la natriurèse aiguë (Koh : $-20$\,\%) tout en restant trop basse en chronique (Lantis $\times 1{,}38$ en moyenne sur 24 h). Deux mécanismes connus ont été ajoutés :
\begin{itemize}
\item \textbf{rénine de la macula densa} ($g_\mathrm{md}$ = 3,76) : le blocage du NKCC2 libère la rénine pendant l'exposition, indépendamment du volume. Le freinage de la natriurèse en prises répétées passe alors surtout par le RAAS, comme le conclut Adin 2021 (« RAAS activation mediates braking », résumé) ;
\item \textbf{latence natriférique} ($\tau_n$ = 1,56 h) : l'action génomique de l'aldostérone sur la réabsorption sodée suit $A$ avec retard.
\end{itemize}
Les 16 constantes ont ensuite été calées ensemble par moindres carrés bornés, sur 17 cibles pondérées par leur incertitude (sonde parallèle, \vp{PDOSE}). L'erreur pondérée passe de 48,5 à 14,9. Le freinage tubulaire structurel ne joue plus qu'un rôle mineur : sans lui, l'erreur passe de 16,6 à 21,9.

Un troisième mécanisme a été essayé puis retiré : un plancher saturant de l'effet natriférique (capacité distale limitée). Le calage le ramène à 0, et la saturation seule aggrave la déplétion.

Contrôles d'échec ajoutés : avec une action natriférique instantanée, la natriurèse de Koh tombe à 1,8 mmol/kg ; sans rénine de la macula densa, la natriurèse de J1 chez Pelligand monte à 9,0 mmol/kg et la créatinine d'Adin à $+107$\,\%.

\subsection{Lantis 2011 en texte intégral (1.3.2)}
Le texte intégral (doi:10.2460/ajvr.72.12.1646, fourni par l'utilisateur, lu) précise le protocole et ajoute des mesures. Groupe furosémide seul : 6 chiens de 30 kg, prises à 7 h et 19 h, régime commercial standard. Mesures :
\begin{itemize}
\item aldostérone/créatinine urinaire : 0,37 $\pm$ 0,14 µg/g au départ ; 0,89 $\pm$ 0,23 à J5 ($\times 2{,}41$) ; 0,95 $\pm$ 0,63 à J10. Ce sont deux urines ponctuelles par jour, à 2–3 h et 8–10 h après la prise du matin, mélangées à parts égales ;
\item kaliémie 4,7 $\to$ 4,2 (J5) $\to$ 4,4 mmol/L (J10) ; poids 30,1 $\to$ 28,9 $\to$ 29,6 kg ; natrémie inchangée ; chlorémie en baisse et bicarbonates en hausse ; PAS en baisse d'environ 15 mmHg à J5 ;
\item phosphorémie, marqueur du DFG selon les auteurs : hausse à J5, retour à la base à J10. Créatinine non mesurée ;
\item les auteurs attribuent l'activation du RAAS à l'action directe du furosémide sur la macula densa, indépendante du volume. C'est le mécanisme $g_\mathrm{md}$ de la 1.3.1.
\end{itemize}
Deux corrections de la comparaison. Le scénario passe à 30 kg. L'aldostérone du modèle est lue aux instants des prélèvements (moyenne de $A\,A_K$ sur les 3 h qui précèdent chacun, puis moyenne des deux), et non plus en moyenne sur 24 h. Avec les constantes de la 1.3.1, le modèle donnait $\times 1{,}99$ (J5) et $\times 1{,}78$ (J10), un $\Delta$K de $-0{,}35$ et un VEC de $-25$ mL/kg à J5 : il retenait trop de Na.

Six cibles ont été ajoutées au calage conjoint : aldostérone, $\Delta$K et poids à J5 et J10 (perte de poids assimilée à celle du VEC). La constante de temps de l'aldostérone ($\tau_A$), placeholder D, a été libérée ; elle revient à 2,06 h. Résultat : $\times 2{,}12$ et $\times 1{,}99$, $\Delta$K $-0{,}42$ et $-0{,}41$, VEC $-2{,}9$ et $-2{,}5$\,\%. L'erreur pondérée sur les 23 cibles passe de 14,5 à 11,2, sans dégrader les autres cibles.

\textbf{Reste} : l'aldostérone de J10 décline dans le modèle ($\times 1{,}99$) alors qu'elle se maintient chez les chiens ($\times 2{,}57$, mais avec un ET de 0,63 µg/g). Le modèle n'a pas de mécanisme qui entretienne l'activation au-delà de la déplétion et de l'exposition, par exemple une hyperplasie de la zone glomérulée.

\subsection{Koh 2021 en texte intégral (1.3.3)}
Le texte intégral (PMC8420016, lu) donne les débits horaires (tableaux 4 à 6). Les valeurs sont des moyennes sur l'intervalle qui se termine à l'instant indiqué (1, 2, 4, 6, 8 h). La natriurèse sur 8 h se calcule donc en sommant débit $\times$ durée : IV 30,1 mmol, soit 2,62 mmol/kg ; per os 40,9 mmol, soit 3,55 mmol/kg. Les cibles précédentes (3,06 et $\approx$ 4,56), obtenues par la méthode des trapèzes, étaient surestimées. Autres données : pharmacocinétique (biodisponibilité orale 55\,\%, cohérente avec le module ; pic plasmatique à 1 h) ; K 4,48 $\to$ 3,97 (2 h) $\to$ 4,15 (8 h) en IV, 4,51 $\to$ 3,97 $\to$ 4,05 per os.

Les chiens étaient à jeun depuis 12 h. Leur kaliurèse de base (0,027 mmol/kg/h) vaut le tiers de celle du modèle, calée sur des chiens nourris (Troìa 2018). Dans le modèle, l'augmentation de kaliurèse est proportionnelle à la sécrétion de base. La cible est donc le rapport de la kaliurèse sur 8 h à la base ($\times 2{,}27$ en IV ; la base orale, 0,92 $\pm$ 1,12 mmol/h, est trop dispersée pour servir). Une première version comparait la kaliurèse nette en valeur absolue et concluait à tort à une kaliurèse deux fois trop forte.

Huit cibles ont été ajoutées : diurèse et natriurèse de la 1re heure IV et des 2 premières heures per os, natriurèse per os, kaliurèse IV rapportée à la base, $\Delta$K à 2 h per os. La constante $a_K$ a été libérée. L'erreur pondérée sur les 31 cibles est de 18,3.

\textbf{Arbitrages.}
\begin{itemize}
\item $a_K$ = 0,49 s'écarte de Young 1983 (0,7, résumé). Avec 0,7 imposé et une cible faible sur la créatinine de Lantis, le freinage part en butée sur ses trois bornes et l'aldostérone de Lantis retombe à $\times 1{,}98$ (erreur 24,2). Les données lues (Lantis, Koh) ont été privilégiées.
\item Créatinine chronique : $+37$\,\% chez Lantis à J10, alors que la phosphorémie, marqueur du DFG selon les auteurs, revient à la base à J10, et qu'Hori 2007 ne montre pas de variation cliniquement pertinente (cité par Lantis ; texte non accessible depuis la session). La hausse chronique de la créatinine est probablement surestimée chez le chien sain. Giorgi (MMVD) donne $+29$\,\% contre $+22$\,\%.
\item La réponse IV reste trop concentrée dans la 1re heure (natriurèse $+29$\,\%). La forme du profil vient surtout de la pharmacodynamie du module furosémide.
\end{itemize}

\subsection{MRC : apports fixes, néphron adapté (remarque de l'utilisateur)}
L'utilisateur a relevé qu'en MRC l'urine n'a pas la natrémie du plasma, que la fuite sodée est faible et que le potassium pose problème. La première version faisait une erreur plus profonde. Elle fixait les apports de Na et de K à l'excrétion du patient avec les FE du chien sain, si bien que l'apport « alimentaire » suivait le DFG (10\,\% de la normale au stade 4).

Correction (\vp{renal_model.json} 1.1.0, puis 1.2.0 pour le K) :
\begin{itemize}
\item les apports sont ceux d'un chien sain de même poids, quel que soit le stade ;
\item les FE de base du patient valent apport $/$ charge filtrée et croissent donc comme $1/\text{ratio}_\text{DFG}$ ;
\item un plafond de sécurité de la FE$_\mathrm{K}$ (\vp{fe_k_max} 1,5, placeholder D, non atteint dans les scénarios validés) ;
\item deux leviers alimentaires : \vp{na_intake_factor} (régime rénal pauvre en sel) et \vp{k_intake_factor} (anorexie, supplémentation) ;
\item les sorties incluent la natriurie et la kaliurie (mmol/L).
\end{itemize}

Résultats (20 kg, sans médicament) :
\begin{itemize}
\item stade 4 ($s$ = 0,85) : FE$_\mathrm{Na}$ 0,20 $\to$ 1,66\,\% ($\times 8{,}3$), FE$_\mathrm{K}$ 9,5 $\to$ 79\,\%, pour un DFG de 70 $\to$ 8 mL/min. C'est cohérent avec des FE « environ dix fois plus élevées au stade IV » (étude de 20 chiens) et des FE qui « dépendent de la créatinine » (Selin 2025). Accès : résumés de recherche ; les textes intégraux (vin.com, pub.epsilon.slu.se) étaient bloqués par le réseau de la session.
\item natriurie de base 63 mmol/L à tous les stades (apport / diurèse), loin des 150 du plasma.
\end{itemize}

Sous furosémide 2 mg/kg IV, à 8 h :
\begin{itemize}
\item VEC : $-17{,}9$ mL/kg chez le chien sain, $-8{,}9$ au stade 2, $-4{,}8$ au stade 3, $-2{,}3$ au stade 4 ;
\item $\Delta$K : $-0{,}39$, $-0{,}31$, $-0{,}19$, $-0{,}08$.
\end{itemize}
La fuite sodée suit la charge filtrée, et l'effet du furosémide, piloté par son excrétion urinaire, baisse aussi.

\textbf{Potassium : réserve proportionnelle au DFG (1.2.0).} Dans la 1.1.0, l'adaptation du néphron restant était parfaite : pour 1,5 fois plus de K dans la ration, la kaliémie montait autant chez le chien sain qu'aux stades 2 et 3. Une recherche dédiée a apporté les données canines suivantes (résumés, sauf mention) :
\begin{itemize}
\item \textbf{Bourgoignie 1981} : chiens à rein restant (DFG 57 $\to$ 16 mL/min), kaliémie à jeun et excrétion de K normales. Après 50 mEq de KCl, la kaliémie monte de $+2{,}2$ à $2{,}5$ (contre $+0{,}9$ à $1{,}2$) et 30–37\,\% de la charge est excrétée en 5 h (contre 61–67\,\%). L'excrétion rapportée au DFG suit une pente commune en fonction de la kaliémie : la \emph{réserve} d'excrétion suit le DFG, l'état de base est adapté.
\item \textbf{Schultze 1971} : excrétion par néphron $\times 6$ à J7 après réduction néphronique, indépendante des minéralocorticoïdes.
\item \textbf{Young 1982, 1983, 1984} : excrétion de K selon la kaliémie à aldostérone fixée (pente forte au-dessus de la normale, faible en dessous) ; aldostérone linéaire en K, nulle à 3,29 mmol/L ; effet de l'aldostérone $\propto A^{0{,}7}$.
\item \textbf{Young 1976} : kaliémie $+8$\,\% avec l'aldostérone, $+39$\,\% sans, pour de grandes variations d'apport.
\item \textbf{King 2017} (lu) : pas de hausse du K sous bénazépril en MRC (49 chiens, jusqu'à 2 ans). \textbf{Segev 2010} : hyperkaliémie chez 47\,\% de 152 chiens en MRC ; K 6,5 sous régime rénal, 5,1 après un régime appauvri en K.
\end{itemize}

Le modèle a été restructuré en conséquence (équations ci-dessus) : réserve $\sigma\cdot\mathrm{DFG}\cdot(K - K_b)$, composante kaliémique de l'aldostérone limitée à la sécrétion de K, charge orale de K. Une première version faisait agir l'aldostérone stimulée par le K sur le sodium. La baisse de K sous furosémide effondrait alors l'aldostérone, et la créatinine d'Adin montait de $+134$\,\%. Le « paradoxe de l'aldostérone » justifie la séparation : la même hormone retient le Na quand elle est portée par l'angiotensine II et ne fait que sécréter le K quand elle est portée par l'hyperkaliémie.

Résultats sur 7 jours, sans médicament :
\begin{itemize}
\item 1,5 fois plus de K dans la ration : $\Delta$K $+0{,}09$ (sain), $+0{,}18$ (stade 2), $+0{,}26$ (stade 3), $+0{,}34$ (stade 4) ;
\item inhibition RAAS de 50\,\% et 1,2 fois plus de K : $+0{,}09$ contre $+0{,}37$ au stade 4.
\end{itemize}
La vulnérabilité au K croît désormais avec le stade.

Limites :
\begin{itemize}
\item le pic de K du rein restant reste sous-estimé ($+1{,}39$ contre $+2{,}2$ à $2{,}5$), car le tampon cellulaire est le même qu'au chien normal ;
\item l'hyperkaliémie spontanée de la MRC avancée vient du module stationnaire (K jusqu'à 6,1 au-delà de $s$ = 0,85), pas du rein dynamique ;
\item la valeur calée de $\sigma_\uparrow$ (1,6) est inférieure à celle tirée directement de Young 1982 et de Bourgoignie ($\approx$ 3).
\end{itemize}
Tests : \vp{renal_ckd_adapts_fractional_excretion}, \vp{renal_potassium_load_scales_with_gfr} (avec un contrôle d'échec sans réserve).

\textbf{IEC et spironolactone.} Aucun module médicament n'est ajouté. Chez le chien sain, le bénazépril abaisse l'activité de l'ECA d'environ 90\,\% (22,1 $\to$ 2,5 U/L) sans changer la diurèse, le freinage ni l'aldostérone pendant une perfusion de furosémide (Adin 2021, lu). Les fortes doses d'énalapril ou de bénazépril n'empêchent pas non plus la hausse de l'aldostérone (Lantis 2015, résumé de recherche). Traduire l'inhibition de l'ECA en baisse de $A$ contredirait donc ces deux sources. Le levier \vp{ace_inhibition} (0–1) de l'API agit directement sur $A$ ; ce n'est pas un médicament.

Placeholders (evidence D) :
\begin{itemize}
\item urine excédentaire iso-natrique (150 mmol/L) ;
\item $n_p$ = 2, plancher de diurèse 0,3, \vp{fe_k_max} 1,5 ($\tau_A$ est calé depuis la 1.3.2).
\end{itemize}

Tests :
\begin{itemize}
\item \vp{renal_matches_literature} : Koh IV, Koh per os, Adin (créatinine à 5 h et à 24 h) ;
\item \vp{renal_chronic_furosemide_pelligand} : natriurèse J1 et J3, diurèse J9 ; Koh IV : diurèse de la 1re heure, natriurèse 8 h ;
\item contrôles d'échec : sans autorégulation, la créatinine d'Adin sort de la fourchette ; sans rétrocontrôle du DFG ($g_\mathrm{DFG}$ = 0), elle ne monte pas ; sans tampon intracellulaire, le $\Delta$K à 2 h sort ; sans composante kaliémique de l'aldostérone, le $\Delta$K à 8 h sort ; sans rénine de la macula densa, la natriurèse de J1 sort ; avec une action natriférique instantanée, la natriurèse de Koh sort ; sans eau libre, la diurèse de J9 sort ;
\item \vp{renal_course_round_trip} (API).
\end{itemize}

Sonde de calage reproductible : \vp{cargo run --release -p vetsim-core --example renal_probe} (scénarios Koh, Adin, charge de K, MRC, Lantis, Wong, Giorgi ; surcharges \vp{clé=valeur}).

\section{Effet sur les références du projet}
\begin{itemize}
\item \textbf{Golden} (\vp{tests/golden/cases.json}), régénéré depuis le Python corrigé (\vp{make_golden.py}). 121 valeurs sur 1782 sortent de la tolérance de l'ancien golden : pneumonie (62), Addison (28), BAV 0,8 (21), MRC 0,8 (2, PaCO$_2$), et 8 cas où seule \vp{lv_edp} ou \vp{la_mean} bouge (MMVD 0,5 et 0,8, pimobendane, bêtabloquant, furosémide, chiens de 5 et 40 kg), sensibles au régime périodique. Rust contre nouveau golden : 1862/1862 valeurs conformes.
\item \textbf{Fixtures Python} (M1–M6, conduction, battements, toxicité, PCG) régénérées. \vp{calibration/calibration_run_python.json} régénéré par \vp{reference/python/calibrate.py}, avec l'accord de l'utilisateur. Le dossier \vp{calibration/} est en principe protégé ; \vp{targets.json} n'a pas été touché.
\item \textbf{Deux tests modifiés}, propriété et seuil conservés :
  \begin{itemize}
  \item \vp{isoflurane_depresses_blood_pressure_and_ventilation_dose_dependently} : le palier « profond » passe de 3 \% à 4 \% de vaporisateur. À 3 \%, la MAC cérébrale n'était que de 1,15 et la PaCO$_2$ montait de 2,8 mmHg, ce qui est conforme à la courbe recalée.
  \item \vp{check_rhythm.py} : la fourchette du BAV porte désormais sur la fréquence observée.
  \end{itemize}
\item \vp{severity_and_weight_sweep_stays_finite} (\vp{hemodynamics.rs}) : la liste des décès attendus passe de vide au chien de 3 kg en MMVD 1,0 (dette d'O$_2$). Ce décès vient du régime périodique, qui élève la POG de la MMVD terminale ; il relève du point « à arbitrer » de la section 2.
\end{itemize}

\section{Recherches et sources bloquées}
\begin{itemize}
\item \textbf{Bloqués} :
  \begin{itemize}
  \item avmajournals.avma.org (Mutoh, McNally, Galloway, Lopez, Johnson en texte intégral) ;
  \item ScienceDirect (\emph{Vet Anaesth Analg}), frontiersin.org, J-STAGE ;
  \item confluence.cornell.edu (de Morais \& DiBartola 1991), vtechworks ;
  \item api.crossref.org (403), OpenAlex (429 intermittents) ;
  \item serveur de vérification bibliographique hébergé (proxy 403) ; PDF JCI de Cohen 1964 et Schwartz 1965 (anti-robot PMC, non contourné).
  \end{itemize}
\item \textbf{Non trouvés} : série canine donnant PAM, FC et DC à 10, 20, 30 et 40 \% de perte sur un même protocole ; bradycardie chiffrée vers 30 \% ; bolus de 20 mL/kg en 15 min chez le chien normovolémique ; PaCO$_2$ chiffrée à 1,5 et 2 MAC en ventilation spontanée ; FE/FI à 5, 15 et 30 min chez le chien ; pentes de DiBartola avec leurs tolérances.
\item « Non trouvé » ne veut pas dire « n'existe pas ». Les DOI viennent d'Europe PMC ; ils n'ont pas été revérifiés par Crossref (bloqué).
\end{itemize}

\section*{Références}
\small
\begin{itemize}\setlength\itemsep{0pt}
\item Acevedo-Arcique CM et al. 2014. \emph{PLoS One}. doi:10.1371/journal.pone.0106620 (lu).
\item Ambros B et al. 2018. \emph{Vet Anaesth Analg} 45:452. doi:10.1016/j.vaa.2018.03.004 (résumé).
\item Bexfield NH et al. 2008. \emph{J Vet Intern Med} 22:66. doi:10.1111/j.1939-1676.2007.0035.x (résumé).
\item Cardillo JH et al. 2024. \emph{Front Vet Sci} 11:1467043. doi:10.3389/fvets.2024.1467043 (lu).
\item Chalifoux NV et al. 2023. \emph{Can Vet J}, PMC10150554 (lu).
\item Goggs R et al. 2022. \emph{J Vet Intern Med}. doi:10.1111/jvim.16374 (lu).
\item Hall H et al. 2023. \emph{J Vet Intern Med}. doi:10.1111/jvim.16786 (lu).
\item Haskins S et al. 2005. \emph{Comp Med} 55:156 (résumé).
\item Hezzell MJ et al. 2013. \emph{J Small Anim Pract} 54:318. doi:10.1111/jsap.12079 (résumé).
\item Itami T et al. 2019. \emph{PLoS One} 14:e0213414. doi:10.1371/journal.pone.0213414 (lu).
\item Kato K et al. 2024. \emph{Open Vet J}. doi:10.5455/OVJ.2024.v14.i5.20 (lu).
\item Lozano AJ et al. 2009. \emph{Vet Anaesth Analg} 36:220. doi:10.1111/j.1467-2995.2009.00451.x (résumé).
\item Lopez LA et al. 2009. \emph{Am J Vet Res} 70:1339. doi:10.2460/ajvr.70.11.1339 (résumé).
\item Madias NE, Bossert WH, Adrogué HJ 1984. \emph{J Appl Physiol} 56:1640. doi:10.1152/jappl.1984.56.6.1640 (résumé).
\item Mazzaferro EM, Ford RB 2012. In \emph{Kirk \& Bistner's Handbook}. doi:10.1016/B978-1-4377-0798-4.00001-3 (lu, tableau 1-6).
\item McNally EM et al. 2009. \emph{Am J Vet Res} 70:1333. doi:10.2460/ajvr.70.11.1333 (résumé).
\item Menaut P et al. 2005. \emph{J Vet Cardiol} 7:75. doi:10.1016/j.jvc.2005.07.002 (résumé).
\item Moran-Muñoz R et al. 2017. \emph{Can Vet J} 58:729, PMC5479674 (lu, via WebFetch).
\item Mosing M et al. 2025. \emph{Front Vet Sci}. doi:10.3389/fvets.2025.1545683 (lu).
\item Pascoe PJ et al. 2006. \emph{Vet Anaesth Analg} 33:97. doi:10.1111/j.1467-2995.2005.00236.x (résumé ; pourcentages cités par Acevedo-Arcique).
\item Rastas, Zhao \& Johnson 2025. \emph{PLoS One} 20:e0317997. doi:10.1371/journal.pone.0317997 (lu).
\item Redondo JI et al. 2012. \emph{Vet Rec} 171:374. doi:10.1136/vr.100476 (résumé).
\item Sage et al. 2023. \emph{Vet Rec} 193:e2586. doi:10.1002/vetr.2586 (résumé).
\item Seo D et al. 2025. \emph{J Vet Emerg Crit Care} 35:214. doi:10.1111/vec.13466 (lu).
\item Thrasher TN, Shifflett C 2001. \emph{Am J Physiol} 280:R1642. doi:10.1152/ajpregu.2001.280.6.r1642 (résumé).
\item van Bavel JJA et al. 2023. \emph{Int J Cardiol Heart Vasc} 44:101169. doi:10.1016/j.ijcha.2022.101169 (lu).
\item Viitanen SJ et al. 2017. \emph{J Vet Intern Med}. doi:10.1111/jvim.14631 (lu).
\item Adin D, Atkins C, Wallace G, Klein A 2021. \emph{J Vet Intern Med} 35:1245. doi:10.1111/jvim.16097, PMC8163123 (lu).
\item Koh SK et al. 2021. \emph{BMC Vet Res} 17:295. doi:10.1186/s12917-021-02998-4, PMC8420016 (lu).
\item Lantis AC et al. 2011. \emph{Am J Vet Res} (résumé).
\item Troìa R et al. 2018 (FE urinaires du chien sain, lu).
\item Bourgoignie JJ et al. 1981. \emph{Kidney Int} 20:482. doi:10.1038/ki.1981.165 (résumé) ; 1985, \emph{Miner Electrolyte Metab} 11:150 (résumé).
\item King JN et al. 2017. \emph{J Vet Intern Med} 31:1113. doi:10.1111/jvim.14726, PMC5508345 (lu).
\item Schultze RG et al. 1971. \emph{J Clin Invest} 50:1061. doi:10.1172/JCI106577 (résumé).
\item Segev G et al. 2010. \emph{J Vet Intern Med} 24:546. doi:10.1111/j.1939-1676.2010.0488.x (résumé).
\item Young DB et al. 1976, \emph{Am J Physiol} 231:945 ; 1982, 242:F599 ; 1983, 244:F28 ; 1984, 247:E328 (résumés).
\item Young AA et al. 2022. \emph{J Vet Intern Med}. doi:10.1111/jvim.16434 (lu).
\end{itemize}
Autres sources citées (résumés) : Bailey 1987, Bigby 2017, Davis 2017, Ferasin 2010, Granot 1985, Kudo 2024, Lamb 2010, Loen 2022 (lu), Morales 1999, Muir 2003, Rögels 2021, Shen 1990, Silverstein 2005 (secondaire), Simões 2016, Smith 2017, Solano 2006, Souza 2025 (lu), Steagall 2006, Thrasher 1998, Turunen 2019, Ueyama 2009, Valverde 2004, Vatner 1974, Yokota 2019, Chae 2026, Ko 2024 (lu).
\end{document}
